% 斯特凡·巴拿赫（综述）
% license CCBYSA3
% type Wiki

本文根据 CC-BY-SA 协议转载翻译自维基百科\href{https://en.wikipedia.org/wiki/Stefan_Banach}{相关文章}

\begin{figure}[ht]
\centering
\includegraphics[width=6cm]{./figures/13afb1dae6b98ddb.png}
\caption{} \label{fig_STFbnh_1}
\end{figure}
斯特凡·巴拿赫（波兰语发音：[ˈstɛfan ˈbanax]；1892年3月30日—1945年8月31日）是一位波兰数学家，被普遍认为是20世纪最重要、最具影响力的数学家之一。他是现代泛函分析的奠基人之一，也是利沃夫数学学派的核心成员之一。他的主要著作是1932年出版的《线性算子论》（法语原名 Théorie des opérations linéaires），这本书是关于泛函分析一般理论的第一部专著。

巴拿赫出生于克拉科夫的一个具有戈拉尔（Góral，高地人）血统的家庭。他自幼对数学表现出浓厚兴趣，常在课间休息时解数学题。完成中学教育后，他结识了数学家雨果·施泰恩豪斯，二人于1919年共同创立了波兰数学学会，并后来创办了学术期刊《数学研究》。1920年，他在利沃夫理工学院获得助教职位，1922年升任教授，1924年成为波兰学术院院士。巴拿赫还是利沃夫数学学派的共同创立者之一，该学派汇聚了波兰战间期（1918–1939）最著名的一批数学家。

以巴拿赫命名的数学概念包括：巴拿赫空间、巴拿赫代数、巴拿赫测度、巴拿赫–塔斯基悖论、哈恩–巴拿赫定理、巴拿赫–施泰恩豪斯定理、巴拿赫–马祖尔博弈、巴拿赫–阿劳格鲁定理、巴拿赫–萨克斯性质以及巴拿赫不动点定理等。
\subsection{生平}
\subsubsection{早年生活}
斯特凡·巴拿赫于1892年3月30日出生在克拉科夫的圣拉撒路综合医院，当时该地属于奥匈帝国。他出身于一个戈拉尔族的罗马天主教家庭，随后由父亲为他施洗。巴拿赫的父母分别是斯特凡·格雷切克和卡塔日娜·巴拿赫，二人都来自波德哈莱地区。格雷切克是奥匈帝国军队驻克拉科夫的士兵，而关于巴拿赫的母亲，人们所知甚少。根据洗礼记录，她出生于博鲁夫纳，是一名家庭女佣。

不同寻常的是，斯特凡的姓氏取自母亲，而不是父亲，但他的名字则继承自父亲。由于军队规定禁止格雷切克这样的士兵（当时军衔为列兵）结婚，而母亲又贫困无力抚养孩子，这对男女决定让孩子由亲友抚养。斯特凡生命的最初几年由祖母照料，但当祖母生病后，格雷切克安排他由克拉科夫的弗朗齐丝卡·普沃娃和她的侄女玛丽亚·普哈尔斯卡抚养。年幼的斯特凡视弗朗齐丝卡为养母，而玛丽亚则如同姐姐一般。在早年，巴拿赫由法籍知识分子尤利乌什·米恩（Juliusz Mien）辅导，他是普沃娃家的朋友，曾移居波兰，以摄影与将波兰文学翻译成法语为生。米恩教授巴拿赫法语，并极有可能在他早期的数学启蒙中起到了鼓励作用。

1902年，年仅10岁的斯特凡·巴拿赫进入了克拉科夫的第四中学（又称戈茨中学，IV Gymnasium / Goetz Gymnasium）就读。虽然这所学校以人文学科为主，但巴拿赫与他最好的朋友维托尔德·维乌科什（Witold Wiłkosz，后来也成为数学家）几乎把所有课间与课余时间都用来研究数学题。后来，巴拿赫曾将自己对数学兴趣的萌发归功于学校的数学与物理教师卡米尔·克拉夫特博士。

巴拿赫是个勤奋的学生，但偶尔也会拿到较低的成绩——他在入学第一个学期时希腊语就不及格。后来他曾批评性地回忆说，当时学校的数学教师水平并不高。1910年，18岁的巴拿赫通过中学毕业考试（matura，相当于高中毕业文凭），随后前往利沃夫（今乌克兰利维夫），打算进入利沃夫理工学院学习。他最初选择工程学作为主修领域，因为那时他认为数学领域“已无新东西可发现”。他还曾在克拉科夫的雅盖隆大学兼读一些课程。由于必须自食其力、打工维持学业，直到1914年、22岁那年，他才终于正式通过中学毕业考试。

第一次世界大战爆发后，巴拿赫因左撇子与视力不佳而被免除兵役。当俄军向利沃夫发动攻势时，他返回克拉科夫，并在那里度过了余下的战争岁月。他以在当地学校当家教、在书店工作、以及担任筑路工地领班等方式谋生。那段时间他也旁听了一些雅盖隆大学的课程，其中包括著名波兰数学家斯坦尼斯瓦夫·扎伦巴和卡齐米日·若拉夫斯基的讲座，但关于他那段时期的生活，人们所知甚少。
\subsubsection{施泰恩豪斯的发现}
\begin{figure}[ht]
\centering
\includegraphics[width=6cm]{./figures/55d67f14e22236f9.png}
\caption{奥托·尼科迪姆与斯特凡·巴拿赫纪念长椅（由斯特凡·杜萨雕刻），位于波兰克拉科夫。} \label{fig_STFbnh_2}
\end{figure}
1916年，在克拉科夫的普兰蒂公园中，巴拿赫偶然遇到了当时著名的数学家之一——雨果·施泰恩豪斯教授。根据施泰恩豪斯的回忆，他当时正漫步于公园中，忽然听到有人谈论“勒贝格积分”（Lebesgue integral——当时仍是数学界相当新的概念），于是他出于好奇走近探听。结果，他就这样结识了巴拿赫以及奥托·尼科迪姆。施泰恩豪斯被这位自学成才的年轻数学家深深吸引。这次偶遇促成了他们长久的合作与友谊。不久之后，施泰恩豪斯邀请巴拿赫尝试解决几个他自己研究已久却未能攻克的问题。巴拿赫在一周内就全部解决了这些问题，两人随后发表了他们的第一篇合作论文——《傅里叶级数的平均收敛性》。此后，施泰恩豪斯、巴拿赫和尼科迪姆，以及其他几位克拉科夫的数学家（包括瓦迪斯瓦夫·希莱博金斯基、莱昂·赫维斯泰克、阿尔弗雷德·罗森布拉特和弗沃季米日·斯托热克）共同创建了一个数学学会，该学会后来发展成为波兰数学学会。该学会于1919年4月2日正式成立。同时，也正是通过施泰恩豪斯的介绍，巴拿赫结识了他未来的妻子——卢恰·布劳斯。
\subsubsection{战间期}
\begin{figure}[ht]
\centering
\includegraphics[width=6cm]{./figures/8c1fb8b25f347a55.png}
\caption{苏格兰咖啡馆——许多著名利沃夫数学家聚会的场所。} \label{fig_STFbnh_3}
\end{figure}
施泰恩豪斯将巴拿赫引荐进入学术界，并极大地推动了他的事业发展。1918年波兰重获独立后，巴拿赫获聘为利沃夫理工学院的助理教师。在施泰恩豪斯的支持下，他甚至在未正式从大学毕业的情况下便获得了博士学位。巴拿赫的博士论文由利沃夫的扬二世·卡齐米日大学于1920年接收[21]，并于1922年正式出版[22]。论文中提出了泛函分析的基本思想——一个当时刚刚诞生、并很快发展为数学中独立分支的全新领域。在1920年完成的论文中，巴拿赫以公理化方式定义了如今所称的“巴拿赫空间”[23]。这篇论文在学术界引起广泛讨论，使他于1922年正式成为利沃夫理工学院的教授。起初，他担任安东尼·沃姆尼茨基教授的助理，1927年则获得了自己的教授席位。1924年，他被接纳为波兰科学院院士。与此同时，自1922年起，巴拿赫还担任利沃夫大学第二数学讲座的主任。

年轻而才华横溢的巴拿赫，很快吸引了一大批数学家聚集在他周围。这群人常在著名的“苏格兰咖啡馆”聚会，逐渐形成了后来闻名世界的“利沃夫数学学派”。1929年，该学派创办了自己的学术期刊《数学研究》，主要刊载巴拿赫所专攻的领域——泛函分析的研究成果。大约在同一时期，巴拿赫开始撰写他最著名的著作——关于线性度量空间一般理论的第一部专著。该书于1931年首先以波兰语出版[24]，次年又被翻译成法语，在欧洲学术界获得了广泛的关注[25]。这部书也是由巴拿赫及其学派主持编辑的一系列数学专著中的第一本。1924年6月17日，巴拿赫被选为位于克拉科夫的波兰科学与美术学院的通讯院士。
\subsubsection{第二次世界大战}
\begin{figure}[ht]
\centering
\includegraphics[width=6cm]{./figures/65c7913374ba8a39.png}
\caption{巴拿赫之墓，位于利维夫的利恰基夫公墓（波兰语称 Lwów）。} \label{fig_STFbnh_4}
\end{figure}
在纳粹德国与苏联入侵波兰之后，利沃夫被苏联占领，持续了将近两年。自1939年起，巴拿赫成为乌克兰科学院的通讯院士，并与苏联数学家关系良好。[5] 为了能够保留教席并继续从事学术研究，他被要求承诺学习乌克兰语。[26] 1940年，他被苏联任命为利沃夫市议会议员。1941年，德军在“巴巴罗萨行动”中占领利沃夫后，所有大学均被关闭。巴拿赫与多位同事以及他的儿子一同在鲁道夫·魏格尔教授的斑疹伤寒研究所任职，担任“虱子饲养员”。在魏格尔研究所工作，为许多失业的大学教授及其同僚提供了免于被任意逮捕或遣送至纳粹集中营的保护。

1944年，苏联红军在“利沃夫－桑多梅日攻势”中重新夺回利维夫后，巴拿赫重返大学，协助重建战后的校务。然而，由于苏联当局正在将被吞并的波兰东部领土上的波兰居民强制迁回波兰境内，巴拿赫开始准备离开利沃夫，计划迁往克拉科夫。他已被雅盖隆大学承诺授予教授席位。[5] 当时，他也被视为波兰教育部长的候选人之一。[27]1945年1月，巴拿赫被诊断出患有肺癌，获准留在利沃夫。同年8月31日，他病逝，享年53岁。他的葬礼在利恰基夫公墓举行，数百人前来悼念。[27] 战后，他在利沃夫大学的教授席位由弗拉基米尔·莱维茨基继任。[28]
\subsection{学术贡献}
巴拿赫于1920年完成、1922年发表的博士论文，首次以形式化的公理体系确立了“完备赋范向量空间”的概念，并奠定了泛函分析这一数学分支的基础。在这篇论文中，巴拿赫将这类空间称为“E类空间”；但在他1932年出版的著作《线性算子论》中，他将这一术语改称为“B型空间”。这一命名的变化很可能促成了后来人们以他的名字命名这些空间，即“巴拿赫空间”[29]。后来被称为“巴拿赫空间”的理论，最早可以追溯到匈牙利数学家弗里杰什·里斯于1916年发表的研究成果，以及同时期汉斯·哈恩和诺伯特·维纳的贡献。[21] 在短暂的一段时间里，数学文献中曾将完备赋范线性空间称为“巴拿赫–维纳空间”，这一称呼源自维纳本人提出的术语。然而，由于维纳在该领域的研究较为有限，最终“巴拿赫空间”这一名称被普遍采用。[29]

此外，巴拿赫在论文中还提出了著名的巴拿赫不动点定理，该定理建立在夏尔–埃米尔·皮卡尔早期方法的基础之上。后来，他的学生们（例如在巴拿赫–绍德定理中）以及其他数学家（如布劳威尔、庞加莱和伯克霍夫都对其作出了扩展。该定理不要求空间具有线性结构，适用于任意完备的柯西空间（尤其是任何完备度量空间）。[21]
\begin{figure}[ht]
\centering
\includegraphics[width=8cm]{./figures/27009a54cda38467.png}
\caption{} \label{fig_STFbnh_5}
\end{figure}
哈恩–巴拿赫定理是泛函分析中的基本定理之一。[21]与巴拿赫相关的其他重要定理还包括：
\begin{itemize}
\item 巴拿赫–塔斯基悖论
\item 巴拿赫–施泰恩豪斯定理
\item 巴拿赫–阿劳格鲁定理
\item 巴拿赫–斯通定理
\end{itemize}
\subsection{荣誉与纪念}
1946年，波兰数学学会设立了斯特凡·巴拿赫奖（波兰语：Nagroda im. Stefana Banacha），以纪念他的卓越贡献。1992年，波兰科学院数学研究所创立了斯特凡·巴拿赫奖章**，专门授予在数学科学领域取得杰出成就的学者。[30] 自2009年以来，波兰数学学会还设立了国际斯特凡·巴拿赫奖，授予在数学科学领域完成最佳博士论文的数学家，旨在“鼓励并资助最有潜力的青年研究者”。[31]

斯特凡·巴拿赫的名字也被用于多所学校与街道命名，这些地方包括华沙、利维夫、希维尼察、托伦和亚罗斯瓦夫等地。

2001年，天文学家保罗·康巴于1997年发现的一颗小行星被命名为16856 Banach，以纪念他。[32]

2012年，波兰国家银行为纪念巴拿赫的成就，发行了一套纪念币，由罗伯特·科托维奇设计，包括200兹罗提金币、10兹罗提银币及2兹罗提北欧金币。[33]

2016年，为纪念巴拿赫与奥托·尼科迪姆首次与雨果·施泰恩豪斯相遇100周年，波兰克拉科夫普兰蒂公园揭幕了一座纪念长椅雕塑，象征这次偶遇对巴拿赫科学事业的决定性影响。[34]

2019年，波兰昆虫学家马尔钦·卡明斯基以“Machleida banachi”命名了一种新发现的南非黑甲虫属昆虫，以纪念巴拿赫对科学的杰出贡献。[35][36]

2021年，波兰纪录片系列《天才与梦想家》在TVP1与TVP Dokument频道播出，其中一集专门讲述了斯特凡·巴拿赫的生平与成就。[37]

2022年，谷歌涂鸦为纪念巴拿赫获得教授头衔100周年，推出了特别纪念图标。[38]
\subsection{名言}
\begin{figure}[ht]
\centering
\includegraphics[width=6cm]{./figures/47b48b062f6319b2.png}
\caption{} \label{fig_STFbnh_6}
\end{figure}
利沃夫数学学派的另一位数学家斯坦尼斯瓦夫·乌拉姆在其自传中引用了巴拿赫的一句话：

“优秀的数学家能够发现定理或理论之间的类比，而最卓越的数学家能发现类比之间的类比。”[39]

雨果·施泰恩豪斯曾评价道：

“巴拿赫是我最伟大的科学发现。”[40]
\subsection{参见}
\begin{itemize}
\item 封闭值域定理
\item 国际斯特凡·巴拿赫奖
\item 波兰人物列表
\item 波兰数学家名单
\item 以斯特凡·巴拿赫命名的事物列表
\item 波兰科学与技术发展年表
\end{itemize}
\subsection{参考文献}
\subsubsection{引用资料}
\begin{enumerate}
\item  “斯特凡·巴拿赫——波兰数学家”。britannica.com，2023年8月27日。
* Pitici 2019，第23页。
* Chemla, Chorla & Rabouin 2016，第224、225、237页。
* “斯特凡·巴拿赫主页”。kielich.amu.edu.pl。检索于2017年8月19日。
* O'Connor, John J.; Robertson, Edmund F., “Stefan Banach”，《圣安德鲁斯大学数学史档案》（MacTutor History of Mathematics Archive, University of St Andrews）。
* Stachura 1999，第51页。
* Waksmundzka-Hajnos 2006，第16页。
* Duda 2009，第29页。
* Kałuża 1996，第2–4页。
* Kałuża 1996，第1–3页。
* Kałuża 1996，第3页。
* Kałuża 1996，第137页。
* Jakimowicz & Miranowicz 2007，第4页。
* Kałuża 1996，第3–4页。
* Jakimowicz & Miranowicz 2007，第5页。
* Kałuża 1996，第13页。
* Kałuża 1996，第16页。
* Jakimowicz & Miranowicz 2007，第6页。
* Ciesielska & Maligranda 2019，第57–108页。
* Kałuża 1996，第23页。
* Jahnke 2003，第402页。
* Stefan Banach（1922）。《关于抽象集上的运算及其在积分方程中的应用》（*Sur les opérations dans les ensembles abstraits et leur application aux équations integrals*），*Fundamenta Mathematicae*（法语与波兰语）。卷3。
* “克拉科夫与利沃夫之间的数学之桥：斯特凡·巴拿赫如何成为世纪最伟大的数学家之一”。krakow1.one，2022年11月16日。
* Stefan Banach：《线性算子理论》（*Teoria operacji liniowych*）。
* Stefan Banach：《Théorie des opérations linéaires》（法语，《线性算子论》）。
* Urbanek 2002。
* James 2003，第384页。
* Zusmanovich, Pasha（2024）。《数学家东行》（*Mathematicians going east*）。*人文数学杂志*（*Journal of Humanistic Mathematics*），第14卷第1期，第114–167页。doi:10.5642/jhummath.EGCM7534。MR 4758663。见第141–142页。
* MacCluer 2008，第6页。
* 波兰科学院数学研究所：斯特凡·巴拿赫奖章。
* “首位国际巴拿赫奖得主”（波兰语）。存档于2015年6月17日。检索于2022年8月11日。
* Urbanek, Mariusz（2014）。《天才们：利沃夫数学学派》（*Genialni. Lwowska Szkoła Matematyczna*）。华沙：Iskry出版社。ISBN 978-83-244-0381-3。
* “波兰国家银行——纪念币”（PDF，波兰语）。检索于2022年8月11日。
* “克拉科夫普兰蒂公园上最著名数学对话一百周年纪念”（波兰语）。存档于2016年10月19日。检索于2022年8月11日。
* Kamiński, Marcin J.; Kanda, Kojun; Smith, Aaron D.（2019年12月10日）。《Machleida属的分类修订》（*Taxonomic revision of the genus Machleida Fåhraeus, 1870*）。*ZooKeys*（第898卷）：第83–102页。Bibcode:2019ZooK..898...83K。doi:10.3897/zookeys.898.46465。ISSN 1313-2970。PMC 6926426。PMID 31875088。
* “甲虫狂热：在博物馆藏品中发现的被遗忘甲虫，以波兰数学家命名”。*Science in Poland*。检索于2025年8月6日。
* “《天才与梦想家》——TVP1新纪录片系列”（波兰语）。检索于2022年8月11日。
* “斯特凡·巴拿赫：Google Doodle 向波兰数学家致敬”。2022年7月22日。检索于2022年8月11日。
* Ulam, Stanislaw M.（1976）。《一个数学家的冒险》（*Adventures of a Mathematician*）。Scribner出版社，第203页。ISBN 9780684143910。
* Strick, Heinz Klaus（2016年11月3日）。《斯特凡·巴拿赫（1892年3月30日–1945年8月8日）——欧洲的数学》（*Mathematics in Europe*）。译者：David Kramer。欧洲数学学会。存档于2016年11月3日。检索于2024年1月19日。
\end{enumerate}
